
\section{何谓弗洛尔（上）同调}
\label{sec:1-1}

一般而言，弗洛尔同调\liftIndex{I014}是定义在无限（$\infty$）维空间上的半无限（$\infty/2$）维同调理论。它已经在多种情形下得到定义，并已证明是研究辛几何、规范场论和低维拓扑等不同领域中诸多问题的一种极其深刻而有用的工具。

弗洛尔在辛几何中引入弗洛尔同调的主要动机 [Flo88IV]，是证明哈密顿微分同胚（或正合辛微分同胚）不动点的阿诺尔德猜想\liftIndex{I001} [Arn65]。阿诺尔德猜想的一个版本断言：紧流形 $M$ 上的任意哈密顿微分同胚
$\phi\colon(M,\omega)\to(M,\omega)$ 均满足
\begin{equation}
  \#\operatorname{Fix}\phi \geq \sum_k \operatorname{rank} H_k(M),
  \tag{1.1}\label{eq:1-1}
\end{equation}
前提是它的所有不动点均非退化。哈密顿微分同胚（或正合辛微分同胚）是由某个（依赖时间的）哈密顿流的时一映射得到的辛微分同胚。

考虑对角子流形
$\Delta\subset(M\times M,\omega\oplus-\omega)$；它是该乘积中的拉格朗日子流形\liftIndex{I020}。这样也可以转而考察这个问题的交点论版本，而这正是弗洛尔在 [Flo88IV] 中最初采用的进路。

设 $L_1,L_0\subset(M,\omega)$ 是两个紧拉格朗日子流形。对于弗洛尔同调，人们或许会期待如下理想性质。

\par\smallskip\noindent\textup{(1.2.1)}\quad
对于每一对拉格朗日子流形 $(L_1,L_0)$，可以赋予一个（分次）阿贝尔群（或 $\mathbb Z_2$-向量空间），记为 $HF(L_1,L_0)$；我们称 $HF(L_1,L_0)$ 为弗洛尔同调。

\par\smallskip\noindent\textup{(1.2.2)}\quad
弗洛尔同调在哈密顿微分同胚下不变。也就是说，若 $\phi_i$（$i=0,1$）是两个哈密顿微分同胚，则
\[
  HF_*(L_1,L_0)\simeq HF_*(\phi_1L_1,\phi_0L_0).
\]

\par\smallskip\noindent\textup{(1.2.3)}\quad
若 $L_0=L_1=L$，则弗洛尔同调群与拉格朗日子流形 $L$ 的通常同调群相同。也就是说，
$HF_*(L,L)\simeq H_*(L;\mathbb Z_2)$。

\par\smallskip\noindent\textup{(1.2.4)}\quad
若 $L_0$ 与 $L_1$ 横截，则弗洛尔同调 $HF_*(L_1,L_0)$ 是一个链复形的同调群；作为群，该链复形由
\begin{equation}
  CF_*(L_1,L_0)=\bigoplus_{p\in L_1\cap L_0}\mathbb Z_2[p]
  \tag{1.3}\label{eq:1-3}
\end{equation}
给出。

事实上，弗洛尔把 $HF(L_1,L_0)$ 定义为链复形 \eqref{eq:1-3} 的同调，其中边界算子 $\partial$ 的矩阵元
$\langle\partial p,q\rangle$ 由如下计数定义：计数（带角的）伪全纯圆盘的数目，其边界的两个半边分别落在 $L_0$ 和 $L_1$ 上，渐近极限为
$p,q\in L_1\cap L_0$（更多细节见第 2.3 节）。需要完成的主要任务，是论证这个定义确实成立，并证明算子 $\partial$ 的边界性质 $\partial\partial=0$。

在 [Flo88IV] 中，弗洛尔完成了这一分析，定义了弗洛尔同调群 $HF(L_1,L_0)$，并在
\[
  L_0=\phi(L),\qquad L_1=L,
\]
其中 $L$ 满足 $\pi_2(M,L)=0$、$\phi\colon M\to M$ 为哈密顿微分同胚的情形下，证明了上述性质 (1.2.1)--(1.2.4)。

继弗洛尔之后，第二作者吴龙根（Yong-Geun Oh）在 [Oh93,95I] 中把弗洛尔同调的构造推广到极小马斯洛夫数大于 $2$ 的任意一对单调拉格朗日子流形 $(L_1,L_0)$（其定义见第 2.4.5 小节和定理 2.4.42）。他还证明了 (1.2.2) 与 (1.2.4)，但注意到 (1.2.3) 一般并不成立。

两位作者均采用 $\mathbb Z_2$-系数。

初看起来，这些条件颇为苛刻而且技术性很强；与弗洛尔同调的另一版本——哈密顿微分同胚的弗洛尔同调——以及阿诺尔德猜想 \eqref{eq:1-1} 的发展相比，尤其如此。

具体而言，对于微分同胚的情形，任意紧辛流形 $M$ 上已经证明了如下结果：对于任意哈密顿微分同胚 $\phi\colon M\to M$，都可以赋予一个取适当系数的分次阿贝尔群，即弗洛尔同调 $HF(M;\phi)$，使得
\[
  HF(M;\phi)\cong H(M),
\]
并且
\[
  \sum_k\operatorname{rank}HF_k(M;\phi)\leq \#\operatorname{Fix}(\phi).
\]
这里 $\operatorname{Fix}(\phi)$ 是 $\phi$ 的不动点集。

在 [Flo89I] 中，弗洛尔给出了哈密顿微分同胚的 $HF(M;\phi)$ 的定义，并证明当 $(M,\omega)$ 单调时，
$HF(M;\phi)\cong H(M;\mathbb Z)$。由此，他证明了单调情形下的阿诺尔德猜想。（在这个问题上，还应提到查尔斯·康利（Charles C. Conley）与爱德华·曾德尔（Eduard Zehnder）针对辛环面的早期结果 [CoZe83]，以及弗洛尔本人由 [Flo88IV] 推得的 $\pi_2(M)=0$ 情形。）赫尔穆特·霍费尔（Helmut Hofer）与迪特马尔·萨拉蒙（Dietmar Salamon）[HoSa95]，以及第四作者小野薫（Kaoru Ono）[Ono95]，把这一结果推广到了半正辛流形。一般情形下弗洛尔同调的构造以及阿诺尔德猜想的证明，已经由 [FuOn99II, LiuTi98, Rua99] 在 $\mathbb Q$ 上完成，并由第一作者和第四作者在 [FuOn01] 中于 $\mathbb Z$ 上完成。（[FuOn01] 的细节仍有待写出。）由于 $HF(M;\phi)$ 同构于 $M$ 的同调群 $H(M)$，这个群本身并不给出辛流形 $(M,\omega)$ 的新不变量。然而，弗洛尔同调的定义本身可以用于各种情形下的阿诺尔德猜想。

另一方面，不能期待拉格朗日交弗洛尔同调理论与哈密顿微分同胚的理论完全相同；也就是说，不能期待性质 (1.2.1)--(1.2.4) 一般都成立。

第一，拉格朗日交弗洛尔同调所处理的是一对拉格朗日子流形 $(L_1,L_0)$，但它无法对完全任意的一对拉格朗日子流形定义。

第二，与微分同胚的情形不同，相关模空间一般既不带稳定殆复结构，也不能形变为带有这种结构的空间，因而不具典范定向。在余切丛情形，第二作者曾在 [Oh97II] 中研究拉格朗日交弗洛尔同调的定向问题；他改进了弗洛尔在 [Flo88IV] 中的结果，从而在 $\mathbb Z$-系数上定义弗洛尔同调。一般而言，对附着于拉格朗日子流形的伪全纯圆盘（或带边黎曼曲面）的模空间赋予定向并非总是可能的；这是带边黎曼曲面上的柯西--黎曼算子的一个颇为非平凡的族指标问题。由此便引出了拉格朗日子流形上的 spin 结构条件（见下文定义 1.6）。

第三，即使 $HF(L_1,L_0)$ 已有定义，不变量 $HF(L_1,L_0)$ 的（辛）几何意义也会变得不那么清楚。例如，考虑 $\mathbb C^n$ 中的一对紧拉格朗日子流形 $(L_1,L_0)$（$\mathbb C^n$ 虽非紧，但在格罗莫夫 [Grom85] 的意义下是驯服的，因此仍能以同样方式对紧拉格朗日子流形建立弗洛尔理论）。显然，存在一个具有紧支集的哈密顿同痕 $\phi_t$，使得 $\phi_1(L_1)\cap L_0=\varnothing$。由这个事实以及 (1.2.2)、(1.2.4) 可知：如果能够定义一对拉格朗日子流形 $(L_1,L_0)$ 的弗洛尔同调，并要求它在哈密顿同痕下不变，那么它必然是平凡的。因此乍看之下，对于 $\mathbb C^n$ 中拉格朗日子流形这一重要情形，弗洛尔理论似乎不会给出任何有趣的结论。然而，只要知道弗洛尔理论已有定义，就可以用它证明关于 $\mathbb C^n$ 中拉格朗日子流形的各种结果。例如，可以用弗洛尔理论证明格罗莫夫关于正合拉格朗日子流形不存在性的定理 [Grom85]，还可以证明 $\mathbb C^n$ 中紧拉格朗日嵌入的马斯洛夫类所受的一些拓扑障碍，如 [Oh96I,96II] 所示。此一思路后来又由其他人进一步发展，其中最著者为保罗·比兰（Paul Biran）与凯·齐利巴克（Kai Cieliebak）[BiCi01,02]。

因此，弄清何时能够定义弗洛尔同调、何时不能，是一项重要任务。正如 [Oh93] 所指出的，如果仿照弗洛尔在 [Flo88IV] 中的做法，仅通过计数伪全纯圆盘的数目来定义算子 $\partial$，那么 $\partial$ 一般不会满足边界性质，即 $\partial\partial\neq0$。这一事实在 [Ono96] 中被应用于弗洛尔同调的一个变体。用物理语言说，$\partial\partial\neq0$ 可以看成孤子效应（即伪全纯圆盘的存在）导致的 BRST 对称性破缺。本书副标题中的 ``anomaly''（反常）一词指的正是这种（BRST）对称性破缺。这意味着，拉格朗日子流形弗洛尔同调的数学框架将比哈密顿微分同胚的相应框架更为复杂（也更有趣）。

本书确立的主要结果，是关于定义拉格朗日交弗洛尔同调的障碍理论。我们还会说明 $L$ 的同调与弗洛尔同调 $HF(L,L)$ 之间的关系，并给出弗洛尔同调在研究拉格朗日子流形几何中的各种应用。本书的结果对于严格表述孔采维奇 [Kon95I] 提出的著名同调镜像对称猜想也具有奠基作用：该猜想以弗洛尔上同调的框架为基础，更准确地说，以第一作者在 [Fuk93] 中引入的 $A_\infty$ 范畴为基础。本书给出了在一般辛流形上严格构造 $A_\infty$ 范畴的若干最重要步骤。

近来，物理学家和数学家都发表了大量工作，计算附着于给定拉格朗日子流形的（伪）全纯圆盘数（例如见 [KatLi01, LiSo01, AAMV05]）。与 Gromov--Witten 不变量的闭情形一样，这些计算主要使用环面作用下的不动点局部化。现在已经广为人知的是，全纯圆盘数本身并没有良好定义，因为该数依赖所作的各种选择，例如为了使模空间横截而选取的扰动。因此，要从局部化等方法算出的这些数中提取出一个数学上有良好定义的对象，是一个非常非平凡的问题。在本书中，我们通过构造一种代数结构——滤过 $A_\infty$ 代数——解决了这个问题；其结构常数由伪全纯圆盘数给出，并满足各种关系。我们证明，与给定拉格朗日子流形相联系的这个滤过 $A_\infty$ 代数的同伦等价类，不依赖于其中涉及的各种选择。一般而言，作者尚不十分清楚局部化等方法所得的数究竟如何与本书研究的结构相关，尽管两者必然相关这一点相当明显。

本书工作的代数基础是 $A_\infty$ 代数理论，我们将在下一节作进一步说明。詹姆斯·斯塔谢夫（James D. Stasheff）于 20 世纪 60 年代在同伦论中利用菅原正博（Masahiro Sugawara；日文 菅原正博）[Sug57] 的思想，在 [Sta63] 中引入了 $A_\infty$ 代数的概念。爱德华·威腾（Edward Witten）[Wit86]（20 世纪 80 年代末）与巴顿·茨威巴赫（Barton Zwiebach）[Zwi91,97]（20 世纪 90 年代初）注意到，$A_\infty$ 结构是开弦理论的一个自然框架。紧接着，第一作者受到唐纳森的一次演讲启发，独立地把 $A_\infty$ 结构引入弗洛尔理论。当时，无论数学研究者还是物理学研究者，都没有很好地理解这些工作之间的关系。

我们很高兴看到，过去十年两方面的发展已经彻底改变了这一局面，使带 D-膜的（拓扑）开弦与弗洛尔上同调之间的关系清楚得多。如今，$A_\infty$ 结构在这两个领域中的出现不再显得突兀，而是一种自然现象。与此同时，这种 $A_\infty$ 结构与同伦论中相应结构的关系也变得更加清楚。

20 世纪 90 年代初，在弗洛尔理论中采用 $A_\infty$ 结构或此类同调代数的思想刚刚萌生时，它显得如此陌生，以至于辛几何学家并不欢迎。直到近来，对弗洛尔理论中各种代数结构的研究才在研究者中变得自然而且不可回避。另一方面，研究弦论的理论物理学家对于弗洛尔上同调，尤其对于辛几何中产生的弗洛尔上同调，仍然很少有细致理解。（在卡姆朗·瓦法（Cumrun Vafa）[Vaf99]、迈克尔·道格拉斯（Michael R. Douglas）[Dou02] 等人的工作之后，一些物理学家开始对此更加熟悉。）我们诚挚地希望，本书的出版能够进一步改善这种状况。

我们还想指出，20 世纪 90 年代后半期为处理由非线性椭圆偏微分方程定义的模空间而发展出的基本技术（[FuOn99II, LiTi98, Rua99, Sie96]），使人们有可能落实从模空间产生的各种代数结构之系统构造所需的数学细节。特别是，采用仓西结构\liftIndex{I037}与多值截面（或其等价替代物）的方法，对于这样的构造至关重要。仓西结构与多值截面的工具体系由第一作者和第四作者于 1996 年发展出来，目的是处理在各种情形中构造适当的（虚）基本闭链所需的诸多横截性问题。本书以更加非平凡的方式运用这些工具来定义（虚）基本链。

本书是一部研究专著。除非另有说明，其中所有结果都是新的。
